In this paper we present a first feasibility study of accelerating MC generation using FPGAs. We will stablish a comparison for a given physics process on the computing speed and power consumption between CPU, GPU and FGPA platforms. 

The following hardware platforms are considered: 
\begin{itemize}
    \item An \textbf{AMD Versal AI Core series device}, which combines traditional FPGA fabric with an array of AI Engine cores on a single chip. The Versal’s programmable logic is utilized for custom datapaths, while the AI Engines (tightly-coupled VLIW vector processors) are exploited for parallel numeric computations. The design was developed and tested on a Versal development board (VCK190).
\end{itemize}

The performance of these devices will be compared to GPU/CPU implementations that already exist within the MG5aMC project. 

\subsection{Target Hardware Platform}
We are using the Xilinx Versal AI Core Series VCK190 Evaluation Kit, which represents the inaugural evaluation platform for the Versal AI Core family, around the XCVC1902 Adaptive Compute Acceleration Platform (ACAP) device. 

The device features a dual-core ARM Cortex-A72 application processor capable of reaching 1.7 GHz, providing high single-threaded performance for system control and compute-intensive tasks. Complementing this is a dual-core ARM Cortex-R5F real-time processor with functional safety certification options, enabling deterministic, low-latency control for demanding real-time applications. 

The VC1902 device contains 400 AI Engine cores operating at 1 GHz or higher, delivering a peak compute throughput of 133 INT8 TOPs (tera-operations per second). These AI Engines are specialized processors designed specifically for machine learning inference and advanced signal processing, organized in a tightly-coupled array with deterministic, low-latency data movement and configurable memory hierarchies that avoid traditional cache miss penalties. Each AI Engine tile comprises: 
\begin{itemize}
    \item A 512-bit SIMD (Single Instruction Multiple Data) vector processing unit supporting both fixed-point and floating-point arithmetic,
    \item A 32-bit scalar RISC processor for control flow,
    \item 32 KB of local data memory with access to 128 KB of addressable shared memory through neighbouring tile connections,
    \item Dual 256-bit load units and single store units with individual address generators
    \item 16 KB of dedicated program memory-
\end{itemize}

The 1,968 DSP58 Engines provide complementary computational capacity optimized for high-precision floating-point (FP32/FP16) and fixed-point (INT8/INT16/INT24) arithmetic, delivering 2.8 TFLOP/s peak floating-point throughput and 11.8 TOP/s INT8 performance. These DSP engines are distributed across the programmable logic in dedicated DSP columns, enabling simultaneous deployment of DSP and general-purpose logic kernels.

The VCK190 provides 1,968,000 system logic cells, 899,840 LUTs (lookup tables), and extensive embedded memory: 158 MB UltraRAM and Block RAM distributed across dedicated memory columns; and 
4 MB Accelerator RAM shareable across compute engines for ultra-low-latency tensor storage. 

Additionally the device incorporates high-speed transceiver arrays and comprehensive memory capabilities. 

\subsection{Event Implementation}
The typical workflow of any matrix element generator, and MG5aMC in particular implies the following steps. The event's phase space is sampled using a pseudo-random number generator to generate the particles' momenta. Then, for each momenta, the matrix element, which is the probability of the event for the given process, is calculated. Traditionally, these steps occur sequentially, but can, however be calculated in parallel for many events at a time~\cite{Wettersten:2023ekm}. 

The Matrix Element computation is, by far, the most computationally intensive bottleneck. As it has been demonstrated~\cite{Valassi:2025xfn}, it is a highly-parallelizable process and thus resulting in significant speeds-up of the computation. In this work we will explore the feasiblity of achieving such acceleration using ACAP devices. 

\subsubsection{Matrix Element computation}

\subsubsection{Color algebra calculation}

\subsection{Test Processes and Workflows}

\subsection{Benchmark Metrics and Tools}
